\chapter{Howtos}

\section{Mounting Directories Using sshfs}

By  making use  of the  user file  system service  FUSE one  can  easily mount
dirctories from machines which are reachable  via SSH. The only thing to watch
out for, is to add the corresponding user to the fuse user group.

\subsection{Mounting HOME directory of laptop on a cluster node using sshfs}

Due  to the  fact,  that everybody  at  Klagenfurt uses  a  laptop to  perform
programming work,  most people  have local working  copies. In order  to avoid
copying of  these working  copies, it is  most convenient  to mount them  to a
directory  on  the cluster  servers  and to  do  compilation  to local  binary
directories on these machines.

The simplest  way of accomplishing this task  is to use sshfs.  By issuing the
follwoing command  on e.g. on sedici the  home directory of the  user simon on
the laptop mac will get mounted to the directory home\_mac on sedici.

\begin{verbatim}
sedici$ sshfs simon@mac: home_mac
\end{verbatim}

Please  make sure  to  do  the compilation  in  a local  directory  and not  a
subdirectory of home\_mac because not  all file system requirements are met to
do compilation there.

\subsection{Bridging Firewalls using SSH Tunnels}

If a  machine is not  directly reachable it  may become necessary to  bridge a
connection using  cascaded SSH Tunnels. In  the following scenario  we want to
mount the home directory of user  simon on mac at Klagenfurt (which is located
in a  VPN) on the machine  nc2 at Simetris GmbH  (which is also in  a VPN). In
order to do so we use the  intermediate machine lse81 at the LSE department in
Erlangen. The first  step is to connect  to lse81 from mac and  creating a SSH
tunnel which forwards the SSH port 22 on mac to port 22022 on lse81.

\begin{verbatim}
mac$ ssh -R 22022:mac:22 strieben@lse81.e-technik.uni-erlangen.de
\end{verbatim}

The next  step is  to forward the  port 22022  on lse81 to  port 22022  nc2 at
Simetris by issuing the following command on nc2.

\begin{verbatim}
nc2$ ssh -L 22022:localhost:22022 strieben@lse81.e-technik.uni-erlangen.de
\end{verbatim}

In a different shell  on nc2 one can now mount the home  directory on mac to a
directory on nc2 by issuing

\begin{verbatim}
nc2$ sshfs -p 22022 simon@localhost: home_mac
\end{verbatim}

\subsection{Unmounting a Share}

To unmount a sshfs share enter

\begin{verbatim}
$ fusermount -u home_mac
\end{verbatim}

\subsection{Mounting via SSH on other operating systems}
\begin{description}
\item[Linux]    For   most   distributions    there   are    binary   packages
  available.    Otherwise    SSHFS     can    be    compiled    from    source
  \url{http://fuse.sourceforge.net/sshfs.html}

\item[Mac  OS X] On  the Mac  one can  either install  MacFUSE and  SSHFS from
\url{http://code.google.com/p/macfuse/wiki/MACFUSE_FS_SSHFS} or the commercial
(40\$)  solution ExpanDrive  \url{http://www.expandrive.com/mac}. Also  have a
look at Macfusion \url{http://macfusionapp.org/}.

\item[Windows]  On Windows  there exist  two  possible solutions.  There is  a
  similar     project     like     FUSE,     which     is     called     Dokan
  \url{http://dokan-dev.net/en}  and which is  freeware. It  is however  in an
  early  stage of  development  and  maybe very  slow.  The other  application
  ExpanDrive is commercial (40\$) \url{http://www.expandrive.com/windows}.
\end{description}

\subsection{Other FUSE File Systems}

\begin{description}
\item[CurlFtpFS] is for locally  mounting FTP directories 
  \url{http://curlftpfs.sourceforge.net/}

\item[NTFS-3G] read/write access to NTFS file systems on the Mac
  \url{http://www.tuxera.com/products/ntfs-open-source/}.

\item[davfs2]        is        for        mounting        WebDAV        shares
  \url{http://savannah.nongnu.org/projects/davfs2}.

\item[VMDK]  by  using  FUSE   the  vmware-mount  utility  from  Virtual  Disk
  Development Kit \url{http://www.vmware.com/support/developer/vddk/} on Linux
  can mount .vmdk hard disk images. The vmware-mount utility is also available
  on Windows.
\end{description}


\section{Debugging using Eclipse, Totalview, gdb}

\subsection{gdb}
\url{http://www.unknownroad.com/rtfm/gdbtut/gdbtoc.html}

\url{http://heather.cs.ucdavis.edu/~matloff/UnixAndC/CLanguage/Debug.html}

\url{http://www.cs.umd.edu/~srhuang/teaching/cmsc212/gdb-tutorial-handout.pdf}

\url{http://de.opensuse.org/SDB:Debuggen_mit_der_Glibc}

\subsection{Eclipse}
\url{http://download.eclipse.org/tools/cdt/docs/tutorials/debug_tutorial/cdt_w_debug.htm}

\url{http://www.users.fh-sbg.ac.at/~rgraf/downloads/SWE1_LB/Debugging%20in%20C%20und%20C++%20Programmen%20unter%20Eclipse.pdf}

\url{http://informatik.unibas.ch/lehre/ws06/cs106/_Downloads/EclipseDebuggingMM.pdf}

\url{http://www.youtube.com/watch?v=azInZkPP56Q}

\subsection{TotalView}
\url{http://www.roguewave.com/documentation/index.php}

\url{https://computing.llnl.gov/tutorials/totalview/}

\url{http://www.lrz.de/services/software/programmierung/totalview/kurs/Totalview_Fortran.pdf}

\section{CFS++ Coding Style}

There is no very strict coding style for CFS++. The main points which have been agreed upon are
\begin{itemize}
  \item No tabs only spaces.
  \item The indentation width is two characters.
  \item Source encoding is UTF-8.
\end{itemize}

\todo[inline]{Put informations from \url{https://lse17.e-technik.uni-erlangen.de:2001/trac/wiki/StyleGuide} here.}

\subsection{Settings for Emacs}

The following attached shell archive
\attachfile[icon=PushPin,description={Emacs settings.},mimetype=text/plain]{attachments/emacs_settings.txt}
implements the CFS++ policies for Emacs. It also contains modes for editing CMake files, gnuplot files and ANSYS APDL input scripts. It also features the CUA mode, which contains keyboard shortcuts which are more consistent with oder editors (e.g. Ctrl+C for copying). Some changes to the Lisp scripts may be necessary to get them working with XEmacs.

\subsection{Settings for Eclipse}

The necessary settings in Eclipse can be made by accessing Window $\rightarrow$ Preferences $\rightarrow$ General $\rightarrow$ Editors and Workspace (c.f Fig. \ref{fig:eclipse_cfs_settings}).
\begin{figure}[htb] 
  \centering 
  \subfloat[Editor settings.]{
    \label{fig:eclipse_editor_settings}
    \begin{overpic}[height=5cm]{pics/eclipse_editor_settings}
%      \put(10,20){\textcolor{red}{\bf Test}}
    \end{overpic}
    }
  \qquad
  \subfloat[Workspace settings.]{
    \label{fig:eclipse_workspace_settings}
    \begin{overpic}[height=5cm]{pics/eclipse_workspace_settings}
    \end{overpic}
  } 
  \caption{Settings for CFS++ in Eclipse.}
  \label{fig:eclipse_cfs_settings}
\end{figure}

For convenience the corresponding settings files are attached here
\attachfile[icon=PushPin,description={Eclipse editor settings.},mimetype=text/plain]{attachments/org.eclipse.ui.editors.prefs.txt}
\attachfile[icon=PushPin,description={Eclipse workspace settings.},mimetype=text/plain]{attachments/org.eclipse.core.resources.prefs.txt}. The have to copied to the workspace/.metadata/.plugins/org.eclipse.core.runtime/.settings directory and must be renamed to \*.prefs.


\section{First Time Implementing}
\label{sec:first_time_implementing}
This chapter is for first time coders, who want to get acquainted with the
structure of CFS++, how the Finite Element method is structured into C++
classes. We start off with the structure itself and the idea behind it and then
by example explain how to start off implementing. But first some general hints
for starting to implement:
\begin{itemize}
  \item First start off by getting to know the helper programs: svn (see
    \ref{sec:svn}), cmake (see \ref{sec:cmake}), etc. 
  \item Do not be scared, just start off implementing. You can not really
    destroy CFS++ since we use subversion. This is the only way to learn,
    talking to CFS++ coders will not help much since talking about the
    procedures in the code and really seeing how it is done are two separate
    things.
  \item The code is kept in relative good shape, lets keep it that way. Keep
    to the object oriented idea and write well readable code.
  \item CFS++ is big, you can not know everything, so think abstract. If you
    want to implement a new time stepping do not worry on how the mass-,
    stiffness matrix or other components are calculated. You just need to know
    that they are there and how to get them.
  \item For your first implementation clear your schedule at least for a week.
    Do not do other stuff on the side.
  \item Keep a sophisticated coder around you, who knows CFS++. Buy him beer and
    be nice to him.
\end{itemize}
\subsection{The CFS++ Structure}
If you ask a CFS++ coder how the code is structured he will tell you about
forms, pde, drivers etc. For the first time the names are often confusing and
everybody has their own idea of how to structure and name method. But
(allegedly) a lot of thought has been put into it so lets look at it by
regarding the Finite Element method.\\
So lets look at an equation and have a look were which part is being handled.
For that we take a pseudo equation 
\begin{equation}
  \int\limits_\Omega \frac{d}{dt}u\psi dx 
  +\int\limits_\Omega\nabla u\cdot\psi dx =
  \int\limits_\Omega f\psi dx
  \label{eq:pseudoFe}
\end{equation}
with the unknown $u$, the test function $\psi$ and the right hand side $f$. This
should cover the most import features, like mass-, stiffness matrix and time
derivative. So lets look at each part in detail, how it is implemented and were
it can be found. Looking at it discrete it would look like
\begin{equation}
  \int\limits_\Omega \frac{1}{\delta t}u_i^{n+1}\varphi_i\psi_j dx 
  -\int\limits_\Omega \frac{1}{\delta t}u_i^n\varphi_i\psi_j dx 
  +\int\limits_\Omega\nabla u_i^{n+1}\varphi_i\cdot\psi_j dx =
  \int\limits_\Omega f_i^{n+1}\psi_j dx\qquad\forall i,j
  \label{eq:pseudoFeDisc}
\end{equation}
for a simple explicit Euler scheme to discretise the time step. $n$ are the time
steps, with time step size $\delta t$ and $i$ are the discrete nodes.
In matrix vector notation it would be
\begin{equation}
  \frac{1}{\delta t}\mathbf{M}\mathbf{u}^{n+1} + \mathbf{Ku}^{n+1}
  =\frac{1}{\delta t}\mathbf{Mu}^n + \mathbf{F}
  \label{eq:pseudoFeMtx}
\end{equation}
Figure~\ref{fig:cfs_calcStructure} schematically shows the most import classes and
methods which are needed to calculate the PDE. There are more calculations
involve but for sake of abstraction this covers the main features.
\begin{figure}[htb] 
  \centering
  \begin{overpic}[width=1.0\textwidth]{pics/cfs_calcStructure}
%    \put(10,20){\textcolor{white}{\bf The electrostatic man!}}
  \end{overpic}
  \caption{Main structure of how a PDE is implemented and which classes are
  necessary.}
  \label{fig:cfs_calcStructure}
\end{figure}%
Registration and setting up the PDE are done in different routines and it should
be regarded separately. The reason is to only do one thing either to read in
values or to use values but not worry about how they got there, otherwise it
will get to complex and confusing. This interface between the xml-File and CFS++
is sketched in Fig.~\ref{fig:cfs_setPdeStructure}.
\begin{figure}[htb] 
  \centering
  \begin{overpic}[width=1.0\textwidth]{pics/cfs_setPdeStructure}
%    \put(10,20){\textcolor{white}{\bf The electrostatic man!}}
  \end{overpic}
  \caption{Main structure of where data from the xml-file is read to use in the
  PDE.}
  \label{fig:cfs_setPdeStructure}
\end{figure}%

\subsubsection{Forms}
\label{sssec:forms}
To calculate the entries of the matrices in \eqref{eq:pseudoFeMtx} the bilinear
forms need to be calculated as
\begin{equation}
  \int\limits_{\Omega_e} \varphi_i\cdot\psi_j dx
  \label{eq:bilinForm}
\end{equation}
at element $e$. This is an entry of the mass matrix entry of
\eqref{eq:pseudoFeDisc}. The time step size $\delta t$ is handled separately by
the time stepping algorithm (see section \ref{sssec:timeStep}) Form can be found in the directory 
\texttt{CFS++/source/Forms}. These are filled with all different kind of forms
which names already give some information and to which PDE it belongs. As it is
object orientated a \texttt{baseForm.*} can be found with the main methods. Most
import is the \texttt{CalcElemMatrix(\ldots)} method which calculates
\eqref{eq:bilinForm}, the matrix entries from one element.
To implement \eqref{eq:bilinForm} into this method:
\begin{itemize}
  \item get integration points -- the element class will help
  \item loop over all integration points
  \item get test function -- the element class will help
  \item get integration weights -- the element class will help
  \item multiply $\psi_i$ with $\varphi_j$
  \item store result in element matrix
\end{itemize}
\subsubsection{Time Stepping}
\label{sssec:timeStep}
The time stepping is called after the assembly. It calculates
\eqref{eq:pseudoFeMtx} to
\begin{equation}
  \mathbf{Ax}=\mathbf{b}.
  \label{eq:Axb}
\end{equation}
So it needs the time step size and the matrices (in this case mass and
stiffness). It may need some additional solution types, like old time steps,
derivatives, etc. -- these need to be initialised in the constructor. Further
parameters from the xml-simulation file can be gathered (see CFS++/source/PDE/newmark.cc) as
example.
In an \texttt{Init()} step necessary variables will be set. Like the matrix
factors, in the case of \eqref{eq:pseudoFeMtx} the stiffness matrix gets the
factor $1$, mass matrix $0$ and unused matrices like the damping matrix are set
to $0$. See General/environment.hh for a list of matrix types.

Some hints on the methods
\begin{description}
  \item[Predictor] For methods which need a predictor
  \item[Corrector] For methods which need a corrector
  \item[{\color{red} AdvanceTimestep}] After this method the solution vector
    should hold the new solution. The last time step vector should also hold the
    new solution. That way everything is prepared for the next time step. The
    reason is that other methods depend on this guarantee (Restart file writing).
    Basically do:
    \begin{equation}
      x_n := x_{n+1}
    \end{equation}
    Do not do this earlier (e.g. Corrector) this may lead to conflicts with
    Iterative-Coupled problems.
\end{description}
The newmark scheme \texttt{CFS++/source/PDE/newmark.*} has most features and still is easy to
read, so a good start to read yourself into.

\subsubsection{PDE}
\label{sssec:pde}
The PDEs located in \texttt{CFS++/source/PDE} is an interface between user and
the method itself. Example of information that is gathered
\begin{itemize}
  \item Is it plane, 3d, axi \ldots?
  \item Is it linear, nonlinear \ldots ?
  \item Is it a moving mesh \ldots?
  \item what non linear solver should be used?
  \item \ldots
\end{itemize}
Now, not every PDE support them all, and not all PDEs need this information. What
options can be set are given in the xml-Schemes.

Important here is the method \texttt{DefineIntegrators} which gathers all
integrator which are necessary. It gathers the forms, which in our case have
been implemented in section \ref{sssec:forms}. These are added into the assemble
object.\\
One special note. If you have a PDE with multiple physical variable
(Navier-Stokes; pressure and velocity) the assemble object needs to know at
which position in the matrix the bilinear form is to be set. Lets look at
\begin{equation}
  \left(\begin{array}{cc}
    A_{11} & A_{12}\\
    A_{21} & A_{22}
  \end{array}\right)
  \left(\begin{array}{c}
    u_1\\
    u_2
  \end{array}\right)=f
\end{equation}
in which the result is split into two vectors $u_1$ and $u_2$ which are coupled
by the matrix entries $A_{12}$ and $A_{21}$. Each of the sub matrices may be
different bilinear forms and this can be set by
\begin{verbatim}
    actStiffContext_UV->SetResults( results_[#1],results_[#2],actSDList, actSDList);
\end{verbatim}
Were \texttt{\#1} sets the matrix block in vertical direction and \texttt{\#2} in horizontal
direction. \texttt{\hyphenation{actStiffContext}} is a \texttt{BiLinFormContext} object which
gets the form and after further initialisation is then stored into the assemble
object.

\subsection{Examining the code in action}
\label{sec:code_in_action}

The explanations in  the previous section are for a  pseudo equation only. The
best  way to  get  started with  the hands-on  CFS++  experience is  to run  a
debugger  on an  existing  example from  the  test suite,  which most  closely
resembles  the problem  to be  investigated.  To  get information  about which
function  calls another  function, static  callgraphs are  a great  tool. Such
graphs can  be generated  using the callgrind  tool of valgrind.   Valgrind is
primarily a profiling tool for  examining and fixing memory leaks in programs.
By using  its callgrind functionality  one can get out  additional information
about the calling  hierarchy of a program. Let's say we  are interested in the
calling hierarchy of an acoustic problem. We can navigate into a corresponding
directory          within         the         test          suite         tree
(e.g. TESTSUITE/singlefield/acoustics/abc3d)  and start the  CFS++ binary from
within valgrind:

\begin{verbatim}
valgrind --tool=callgrind $CFS_BINARY_DIR/bin/OPENSUSE_11.3_X86_64/cfsbin abc3d
kcachegrind
\end{verbatim}

The first  command starts  the CFS++ process  within valgrind and  records the
complete calling history. The second command starts up a GUI in which the call
graph can be  examined visually as Fig.  \ref{fig:kcachegrind}  shows.  

\begin{figure}[htb] 
  \centering 
  \label{fig:kcachegrind}
  \begin{overpic}[height=8cm]{pics/kcachegrind_acou_define_int_callgraph}
  \end{overpic}
  \caption{Callgraph for AcousticPDE::DefineIntegrators() in KCachegrind.}
\end{figure}

While this
display conveys already much information  about the calling hierarchy, it does
not give too many hints about the transient behaviour of the program. In order
to get that kind of information, it is  a good idea to start a debugger on the
same problem.  In  Klagenfurt we have a license for  the TotalView debugger. A
license for  TotalView is also available  on the RRZE machines  in Erlangen. A
nice feature, which sets TotalView apart  from other debuggers is, that it can
also step backwards in time. This  feature is especially useful for a beginner
in CFS++,  since the process  does not have  to be restarted everytime  to get
back to a previous breakpoint. To  start CFS++ in TotalView, use the following
command:

\begin{verbatim}
totalview $CFS_BINARY_DIR/bin/OPENSUSE_11.3_X86_64/cfsbin -a abc3d
\end{verbatim}

It is important to enable the replay engine feature in the window that pops up
(c.f. Fig. \ref{fig:totalview_enable_replay_engine}).

\begin{figure}[htb] 
  \centering 
  \label{fig:totalview_enable_replay_engine}
  \begin{overpic}[height=6cm]{pics/totalview_enable_replay_engine}
  \end{overpic}
  \caption{Enabling backward stepping in TotalView.}
\end{figure}

Afterwards one can  set breakpoints in the main window  and start the process.
Figures show the main window and  the callgraph window (Tools $->$ Call Graph)
of   TotalView  (c.f.   Fig.   \ref{fig:totalview_acou_define_int_callgraph}).
Please note,  that next  to the  forward stepping buttons  there are  also the
backward      stepping     pendants      available     in      the     toolbar
(Fig. \ref{fig:totalview_acou_define_int}).

\begin{figure}[htb] 
  \centering 
  \subfloat[TotalView stopped at DefineIntegrators() in AcousticPDE.]{
    \label{fig:totalview_acou_define_int}
    \begin{overpic}[height=8cm]{pics/totalview_acou_define_int}
%      \put(10,20){\textcolor{red}{\bf Test}}
    \end{overpic}
    }
  \qquad
  \subfloat[Corresponding call graph.]{
    \label{fig:totalview_acou_define_int_callgraph}
    \begin{overpic}[height=8cm]{pics/totalview_acou_define_int_callgraph}
    \end{overpic}
  } 
  \caption{CFS++ running the abc3d example in TotalView.}
  \label{fig:totalview_abc3d}
\end{figure}



\subsection{First implementation by example}
\subsection{Implementing a new PDE}
\subsection{Implementing a new solver}
\subsection{Implementing a new time stepping scheme}
\subsection{Implementing new readers/writers}
\section{Write Debug Output to Log Files}
\url{https://lse17.e-technik.uni-erlangen.de:2001/trac/wiki/Logging}

\section{Adding new Results}
\section{Video Howtos}
